\documentclass[11pt,a4paper]{article}
\usepackage[left=2.2cm,right=2.2cm,top=2.2cm,bottom=2.2cm]{geometry}
\usepackage{fontspec}
\usepackage[spanish,es-noquoting]{babel}
\usepackage{booktabs}
\usepackage{longtable}
\usepackage{array}
\usepackage{enumitem}
\usepackage{xcolor}
\definecolor{acc}{RGB}{0,90,160}
\definecolor{warn}{RGB}{170,60,20}
\usepackage{titlesec}
\titleformat{\section}{\large\bfseries\color{acc}}{\thesection.}{0.6em}{}
\titleformat{\subsection}{\normalsize\bfseries}{\thesubsection.}{0.5em}{}
\titlespacing{\section}{0pt}{14pt}{6pt}
\usepackage{fancyhdr}
\pagestyle{fancy}\fancyhf{}
\fancyhead[L]{\footnotesize Fagoterapia MDR/XDR/PDR \textit{P.\ aeruginosa} --- estado del proyecto}
\fancyhead[R]{\footnotesize\thepage}
\renewcommand{\headrulewidth}{0.4pt}
\usepackage[hidelinks]{hyperref}
\setlength{\parskip}{5pt}
\setlength{\parindent}{0pt}
\newcommand{\ok}{\textcolor{acc}{\textbf{✓}}}
\newcommand{\no}{\textcolor{warn}{\textbf{✗}}}

\begin{document}

\begin{center}
{\Large\bfseries Fagoterapia para \textit{Pseudomonas aeruginosa} MDR/XDR/PDR}\\[2pt]
{\large Revisión sistemática y metaanálisis de proporciones}\\[6pt]
{\normalsize Estado del proyecto --- 3 de agosto de 2026}\\[2pt]
{\footnotesize Universidad Católica de Cuenca}
\end{center}

\vspace{4pt}
\hrule
\vspace{10pt}

\section{Dónde estamos}

El corpus analítico está cerrado y estable: \textbf{40 brazos de estudio, 33
estudios, 91 pacientes}, todos observacionales. El manuscrito compila sin errores
y las seis compuertas de calidad automáticas pasan.

Esta sesión no tocó el análisis. Se dedicó a \textbf{reconstruir la arquitectura
de búsqueda}, que era el defecto más grave que arrastraba la revisión, y a
resolver una revisión por pares Q1 de doble especialidad.

\section{Arquitectura de búsqueda: ocho fuentes ejecutadas}

Cada fuente se busca con \textbf{dos brazos}. El brazo A exige el organismo; el
brazo B lo omite deliberadamente.

\begin{center}
\begin{tabular}{@{}lrrl@{}}
\toprule
\textbf{Fuente} & \textbf{Brazo A} & \textbf{Brazo B} & \textbf{Aporte no duplicado} \\
\midrule
PubMed / NCBI            & 1\,481 & 8\,532 & 9\,561 (corpus congelado) \\
Scopus                   & 2\,875 & 9\,372 & pendiente deduplicar \\
ProQuest                 & 1\,013 & 3\,403 & pendiente deduplicar \\
Cochrane CENTRAL         &     53 &    304 & 304 \\
ClinicalTrials.gov       &     17 &    121 & 121 \\
BVS (LILACS/CUMED/BINACIS) & --- & n/a & $\sim$454 \\
EudraCT                  & --- & --- & 3 \\
CTIS                     & --- & --- & 5 \\
\bottomrule
\end{tabular}
\end{center}

\subsection{Por qué hacen falta dos brazos}

Cuatro estudios incluidos --- \textbf{Pirnay 2024}, Green 2023, Leitner 2021 y
Liu 2025 --- son cohortes multipatógeno que \textbf{no nombran el organismo} en
título, resumen ni vocabulario controlado. Una consulta que exija ``Pseudomonas
aeruginosa'' no puede alcanzar un artículo titulado \textit{``Personalized
bacteriophage therapy outcomes for 100 consecutive cases''} por mucha
sensibilidad que se le ponga. Medido en PubMed: el brazo A cubre 36 de 40, el
brazo B cubre 29, y \textbf{sólo la unión cubre 40}.

Pirnay es el mayor contribuyente del corpus. Sin el brazo B se pierde.

\subsection{Tres fuentes ausentes, con evidencia}

\begin{itemize}[leftmargin=1.4em,itemsep=2pt]
\item \no~\textbf{Embase} --- no está en la suscripción institucional. El portal
ofrece Ovid MEDLINE, Ovid Español y Books Ovid Odontología; ninguno es Embase.
\emph{Mitigación parcial:} CENTRAL sí lleva registros que Cochrane obtiene de sus
propias búsquedas en Embase --- pero sólo ensayos controlados, no reportes de
caso, que son el 90\,\% de este corpus.
\item \no~\textbf{Web of Science} --- sin vía encontrada. Research4Life lista a
Clarivate, pero como \textit{Current Contents} (servicio de alertas de índices,
sin índice de citas ni exportación para cribado). Quedan sin comprobar la vía
CEDIA y el acceso de terceros.
\item \no~\textbf{EBSCO} --- sólo dos colecciones de libros electrónicos. No hay
CINAHL, ni MEDLINE-vía-EBSCO, ni Academic Search. La ecuación se retiró porque no
tenía dónde ejecutarse.
\end{itemize}

\section{Cribado}

\begin{center}
\begin{tabular}{@{}lr@{}}
\toprule
Corpus congelado de PubMed (denominador PRISMA) & 9\,561 \\
Excluidos en etapa 1 por regla determinista & 2\,402 \\
\quad fago como técnica de laboratorio & 2\,132 \\
\quad tipo de artículo sin datos primarios & 140 \\
\quad indexado animal y no humano & 130 \\
\midrule
\textbf{Pasan a cribado por título/resumen} & \textbf{7\,159} \\
\bottomrule
\end{tabular}
\end{center}

Las reglas se validaron contra los \textbf{42 estudios incluidos conocidos}:
ninguna regla excluye a ninguno, y el script se niega a aplicarse si esa
validación falla.

\textbf{Se probaron cuatro filtros de contenido y tres fallaron.} Exigir MeSH
\textit{Humans} pierde 11 positivos conocidos (los registros recientes aún no
están indexados); exigir palabra clave clínica pierde 5; la disyunción de ambos
todavía pierde 2. Sólo ``fago como técnica de laboratorio'' no pierde ninguno.
Las tres reglas rechazadas quedaron documentadas para que nadie las reintente.

\subsection{Plataforma}

Reseña creada en \textbf{Rayyan} (Ouzzani et al., 2016), en modo ciego, con los
archivos RIS listos para importar. Rayyan aporta doble cribado independiente,
registro nativo de decisiones y generación del diagrama PRISMA.

\section{Revisión por pares Q1 --- veredicto}

Se ejecutó una revisión de doble especialidad (microbiología de \textit{P.\
aeruginosa} + metaanálisis de proporciones). \textbf{Veredicto: Major Revision},
y \textbf{en contra} de reclasificar como scoping review.

Antes de argumentar, se verificaron las cifras del propio informe contra el
pipeline. \textbf{Dos no sobrevivieron:}

\begin{enumerate}[leftmargin=1.5em,itemsep=3pt]
\item El informe decía que el estrato MDR colapsa a ``3 estudios / 4 pacientes''
bajo clasificación verificada. Son \textbf{2 estudios y 3 pacientes} --- el
colapso es peor de lo que el propio referee creía.
\item El informe leía la divergencia entre modelos (FT 0{,}5\,\%; GLMM 9{,}2\,\%;
DL 24{,}0\,\%) como posible señal de que los datos no admiten \emph{ningún}
modelo. La aritmética dice lo contrario: el GLMM reproduce la proporción cruda
7/76 = 9{,}21\,\% \textbf{exactamente}, y DL coincide con el modelo de efecto fijo
porque $\hat\tau^2 = 0$ colapsa uno sobre otro. No son tres estimadores
compitiendo: es una verosimilitud correcta más dos modos de fallo con nombre y
cita. \textbf{La Tabla 6 es evidencia \emph{a favor} del modelo primario.}
\end{enumerate}

\subsection{Los tres cambios críticos}

\begin{enumerate}[leftmargin=1.5em,itemsep=3pt]
\item \textbf{Tabla estructurada ECA vs.\ observacional en el texto principal.}
Los cuatro ensayos aleatorizados son la única evidencia comparativa del campo y
su señal es nula o negativa --- pero están en prosa, mientras la literatura de
uso compasivo, cuya publicación depende del resultado favorable, conserva estatus
cuantitativo. \emph{La evidencia positiva se agrupa y la negativa se narra.}
\item \textbf{Degradar las estimaciones estratificadas por resistencia} a
análisis de sensibilidad, y llevar el colapso verificado al resumen.
\item \textbf{En las 11 celdas con $\hat\tau^2 = 0$}, dejar de llamarlas efectos
aleatorios y suprimir el intervalo de predicción, que hoy se imprime idéntico al
IC e implica una propagación de incertidumbre que no ocurre.
\end{enumerate}

\section{Revista objetivo}

\textbf{\textit{Clinical Microbiology and Infection} (CMI)}, revista oficial de
ESCMID. Elegida sobre JAC por audiencia: los dos hallazgos de cabecera son de
\textbf{medición microbiológica} --- que las etiquetas de Magiorakos no
sobreviven a la verificación independiente, y que la erradicación no es un solo
constructo --- y el lector de CMI asigna esas etiquetas profesionalmente.

Requisitos leídos de la guía de autores de CMI el 3 de agosto de 2026:

\begin{center}
\begin{tabular}{@{}lrr@{}}
\toprule
 & \textbf{CMI} & \textbf{Manuscrito} \\
\midrule
Texto principal (revisión sistemática) & 3\,500 palabras & \textcolor{warn}{\textbf{20\,639}} \\
Abstract estructurado & 300 máximo & 140, sin estructurar \\
Checklist PRISMA 2020 & obligatoria al enviar & pendiente \\
\bottomrule
\end{tabular}
\end{center}

Hay que recortar el \textbf{83\,\%}. Para calibrar: CMI limita los
\textit{Original Articles} a 2\,500 palabras, así que ``Systematic Review'' ya es
el formato más generoso de la revista.

\textbf{Un requisito que muerde:} CMI exige que el abstract estructurado incluya
una sección \textbf{``Interventions''}, y esta revisión \emph{no caracteriza la
intervención} --- no registra producto, dosis, duración ni si fue personalizada.
Es una limitación ya reconocida en Methods, pero ahora ocupa un campo obligatorio.

\section{Lo que falta}

\subsection{Depende del equipo}

\begin{itemize}[leftmargin=1.4em,itemsep=3pt]
\item \textbf{Registro en PROSPERO.} CMI lo pedirá; el ítem 24a de PRISMA exige
el número. Se puede hacer ya.
\item \textbf{Segundo revisor independiente.} Rayyan está en modo ciego, pero el
doble cribado \emph{sólo existe si hay dos personas}. Es la objeción que los
referees marcan desde la ronda 5 y la única que no se cierra con trabajo técnico.
\item \textbf{Exportar las siete fuentes restantes} a Rayyan. Scopus en tandas de
2\,000; en BVS hay que desmarcar MEDLINE antes de exportar o se cuenta dos veces.
\item Anotar las \textbf{seis bases de ProQuest} --- ya identificadas: Coronavirus
Research Database, Ebook Central, Education Research Index, PRISMA Database with
HAPI Index, ProQuest Central y Publicly Available Content Database.
\end{itemize}

\subsection{Trabajo técnico pendiente}

\begin{itemize}[leftmargin=1.4em,itemsep=3pt]
\item Deduplicación cruzada sobre \textbf{DOI, PMID y NCT}. Las tres claves son
necesarias: CENTRAL ingiere registros de ClinicalTrials.gov y BVS incluye
MEDLINE, así que esos solapes son estructurales.
\item Reconciliar las decisiones de Rayyan contra el corpus congelado.
\item Recorte a 3\,500 palabras con reubicación a material suplementario.
\item Reescritura de Methods y del diagrama PRISMA con la arquitectura definitiva.
\end{itemize}

\textbf{Orden recomendado:} no reescribir Methods antes de terminar el cribado.
Los conteos por fuente y por etapa son lo que imprime el diagrama, y si cambian
después habrá que rehacer el texto con riesgo de dejar una cifra vieja.

\section{Advertencias registradas en esta sesión}

Cuatro trampas que producen números equivocados sin dar ninguna señal de error:

\begin{enumerate}[leftmargin=1.5em,itemsep=3pt]
\item \textbf{El traductor de Chrome corrompe consultas.} En Scopus reescribió el
valor real del campo: \texttt{phage*} pasó a \texttt{fago*}, \texttt{AND} a
\texttt{Y}. La consulta se ejecuta y devuelve un número plausible para algo que
nadie escribió. \emph{Solución:} pasar la consulta por URL y verificar que la
página la reimprime en inglés.
\item \textbf{``Contain any of these terms'' de CTIS no es un OR.}
\texttt{bacteriophage phage} devuelve 0; \texttt{bacteriophage} solo devuelve 5.
Añadir un término vació la búsqueda.
\item \textbf{Un cero puede ser un campo vacío.} El primer intento en CTIS
devolvió ``0 results'' porque el texto nunca llegó al campo. Un nulo real y un
fallo de interfaz se ven idénticos. \emph{Regla:} ningún cero vale sin control
positivo (\texttt{cancer} devuelve 3\,036).
\item \textbf{Las etiquetas de pestaña mienten.} ``Colección LILACS Plus'' de BVS
incluye MEDLINE: de sus 2\,379 registros, 1\,815 son MEDLINE y duplicarían el
brazo de PubMed. El desglose real está en la faceta de base de datos.
\end{enumerate}

\vspace{8pt}
\hrule
\vspace{6pt}
{\footnotesize Documentación completa en \texttt{quality\_reports/}:
\texttt{search\_equations\_all\_databases.md} (ecuaciones y conteos por fuente),
\texttt{2026-08-03\_referee\_q1\_micro\_stats.md} (informe de revisión),
\texttt{journal\_targeting.md} y \texttt{cmi\_submission\_requirements.md}
(revista objetivo), \texttt{document\_corpus\_provenance.md} (auditoría de
procedencia).}

\end{document}
